Dear Stefanie, Anders, Shi, Philip, and Lin,
As I mentioned to you before, there is no obligation that you all have to handle articles submitted to the special issue.   But just in case there is interest (and you are not on the author list of the submission), I send to you the current list of submissions.  If after one week there is no interest in given submission(s), then I will handle those  submissions.

List follows ...
Kind regards,
Mark



JCOMP-D-25-01231
Quantum multiple-valued decision diagram; Reversible logic; Quantum circuit
Xing-Qiang Zhao
Sun Yat-Sen University
CHINA
Xing-Qiang Zhao
Nanrun Zhou
Anwei Luo
Tianlong Chen

The most common representation of quantum mechanical phenomena is
transformation matrices. However, the transformation matrices grow exponentially with
the size of a quantum system. which poses significant challenges for efficient
representation and manipulation of quantum functionality. To address this issue,
quantum decision diagrams (QDDs) have been put forward to represent quantum
systems in terms of decision diagrams. In the realm of classical decision diagrams, the
idea of multi-output function description in the form of binary decision diagram (BDD)
enables to carry out a simple analysis of multi-output function by reversible logic, which
is dedicated to single-output functions. However, the existing QDDs lack a formal basis
and do not allow the change of the variable order, an established core functionality in
decision diagrams. Besides, it is extremely complex and difficult to alter the variable
order, which is crucial to determine more compact representations. Therefore the full
potential of quantum multiple-valued decision diagram (QMDDs) or decision diagrams
has not been fully exploited yet for quantum functionality in general. To conveniently
change the variable order, inspired by QMDDs and reversible logic, a novel quantum
multiple-valued decision diagram was designed based on reversible logic, i.e.,
QMRLDD. The formal definitions of QMRLDD is provided, and the QMRLDD is proved
equivalent to quantum circuits, which guarantees its correctness and canonicity.
Moreover, the spatial complexity of QRMLDD is better than that of other typical
QMDDs based on matrices. The feasibility of QMRLDD on quantum computers is
proved also, which provides a new idea for the development of quantum compiler.
Experimental results confirm our theoretical results.

-------------------------------------------------------------

JCOMP-D-25-01184
Bayesian calibration; marginal likelihood; large data; stability; quantum
characterization
Mohammad Motamed
University of New Mexico
Albuquerque, UNITED STATES
Mohammad Motamed
N. Anders Petersson
We propose a marginal likelihood strategy within the Kennedy-O’Hagan (KOH)
Bayesian framework, where a Gaussian process (GP) models the discrepancy
between a physical system and its simulator. Our approach introduces a novel
marginalized likelihood by integrating out the degenerate eigenspace of the covariance
matrix, rather than approximating the original likelihood. Unlike approximation methods
that compromise accuracy for computational efficiency, our method defines an exact
likelihood—distinct from the original but preserving all relevant information. This
formulation achieves computational efficiency and stability, even for large datasets
where the covariance matrix nears degeneracy. Applied to the characterization of a
superconducting quantum device at Lawrence Livermore National Laboratory, the
approach enhances the predictive accuracy of the Lindblad master equations for
modeling Ramsey measurement data by effectively quantifying uncertainties consistent
with the quantum data.

-------------------------------------------------------
JCOMP-D-25-01132
quantum computing, quantum optimal control, hermite method, discrete adjoint,
gradient, optimization
Spencer Lee
Michigan State University
East Lansing, MI UNITED STATES
Spencer Lee
Daniel Appelo, PhD
This work introduces the High-Order Hermite Optimization (HOHO) method, an openloop discrete adjoint method for quantum optimal control. Our method is the first of its
kind to efficiently compute exact (discrete) gradients when using continuous,
parameterized control pulses while solving the forward equations (e.g. Schrodinger's
equation or the Linblad master equation) with an arbitrarily high-order Hermite RungeKutta method. The HOHO method is implemented in QuantumGateDesign.jl, an opensource software package for the Julia programming language, which we use to perform
numerical experiments comparing the method to Juqbox.jl. For realistic model
problems we observe speedups up to 775x.

---------------------------------------------------------------
JCOMP-D-25-01196
Standard Sylvester-conjugate matrix equations, Sampling discretion, Space
compressive approximation, Zeroing
neural dynamics, Euler-forward formula，Probability theory.
Dongqing Wu
Zhongkai University of Agriculture and Engineering
CHINA
Jiakuang He
Dongqing Wu
Time-variant standard Sylvester-conjugate matrix equations are presented as early
time-variant versions of the complex conjugate matrix equations. Current solving
methods include Con-CZND1 and Con-CZND2 models, both of which use ode45 for
continuous model. Given practical computational considerations, discrete these models
is also important. Based on Euler-forward formula discretion, Con-DZND1-2i model
and Con-DZND2-2i model are proposed. Numerical experiments using step sizes of
0.1 and 0.001. The above experiments show that Con-DZND1-2i model and ConDZND2-2i model exhibit different neural dynamics compared to their continuous coun
terparts, such as trajectory correction in Con-DZND2-2i model and the swallowing
phenomenon in Con-DZND1-2i model, with convergence affected by step size. These
experiments highlight the differences between optimizing sampling discretion errors
and space compressive approximation errors in neural dynamics.

--------------------------------------------------
